
\chapter{Discussion}

This chapter interprets the separate evaluations of meaning representation and
surface realization. The results support several practical conclusions, but
they also show that the choice of method depends on the intended objective:
semantic coverage, predicate control, reference similarity, judge-oriented
faithfulness, or inspectability.

\section{Meaning representation}

\subsection{Pipeline comparison}

The pipeline results identify Text-First Extraction with Normalization as the
most reliable of the three evaluated methods. Its automatic additions and
omissions scores are 4.994 and 4.992, respectively, with an overall score of
4.996 in Table~\ref{tab:pipeline-text-first-triples-results}. The same ranking
is retained in the human evaluation: its overall score is 4.913, compared with
4.400 for Evolving Catalog and 4.175 for Blueprint-Iterative Refinement in
Tables~\ref{tab:human-blueprint-results}, \ref{tab:human-text-first-results},
and~\ref{tab:human-catalog-results}. This agreement between the two aggregate
evaluations is useful evidence that the ranking is not solely an artifact of
the automatic judge.

The result is consistent with the design of Text-First Extraction with
Normalization. It allows each instance to be extracted independently, then
normalizes the resulting predicates after extraction. This avoids constraining
an instance to a vocabulary that was created from earlier cases. The method
uses 201.2 unique predicates on average, compared with 44.6 for Evolving
Catalog and 22.0 for Blueprint-Iterative Refinement
(Tables~\ref{tab:pipeline-blueprint-triples-results}, \ref{tab:pipeline-text-first-triples-results},
and \ref{tab:pipeline-catalog-triples-results}). A broader vocabulary is
compatible with the semantic diversity of the input data and is plausibly one
reason why fewer facts are omitted. However, the predicate count is a
descriptive result rather than a controlled ablation; it cannot establish that
vocabulary size alone causes the observed quality difference.

The errors of the other two pipelines illustrate the cost of early vocabulary
control. Evolving Catalog is second overall, but its automatic omission scores
drop to 3.920 for ice hockey and 3.850 for OWID. Its human omission score is
particularly low for OWID at 2.20 and is also reduced for OpenWeather at 3.70
(Tables~\ref{tab:pipeline-catalog-triples-results} and
\ref{tab:human-catalog-results}). The stable catalog provides a coherent and
auditable vocabulary, but facts that do not fit its frozen or slowly expanding
catalog can be lost. Blueprint-Iterative Refinement,
provides the strongest up-front specification and therefore the clearest
human-inspectable artifact. Its omission results nevertheless show the limits
of a compact predetermined blueprint: automatic scores are 3.380 for ice
hockey and 3.840 for OpenWeather, while the human omission mean is 3.825
(Tables~\ref{tab:pipeline-blueprint-triples-results} and
\ref{tab:human-blueprint-results}).

The shared data-to-text stage is close to the ceiling of the automatic scale:
the mean summary and faithfulness scores are 4.996 and 4.974 in
Table~\ref{tab:pipeline-text-results}. The human evaluation is more critical,
with means of 4.50 and 4.20 in Table~\ref{tab:human-text-results}. In
particular, GSMArena and OpenWeather both receive a human faithfulness mean of
3.60. This indicates that their generated summaries contain meaningful
faithfulness issues under the evaluation rubric, but the mean should not be
read as the rating of a single instance. It aggregates ten author-assigned
ratings per domain. The gap between automatic and human means therefore
motivates a cautious interpretation of the apparently near-perfect automatic
scores.

\subsection{Fine-tuning as pipeline distillation}

Fine-tuning substantially improves the model's ability to produce triples from
raw data. The base Gemma model obtains a mean triple F1 of only 0.0205, whereas
the three-epoch baseline reaches 0.2340
(Tables~\ref{tab:finetune-base-results} and
\ref{tab:finetune-baseline-results}). One epoch already improves F1 to 0.1209,
showing that the task-specific targets provide a strong adaptation signal even
with limited training. At the same time, the results do not select a single
configuration as best for every purpose. The three-epoch baseline has the
highest reported average triple F1, the one-epoch model has the highest
predicate adherence, and the low-learning-rate configuration has the highest
average text BLEU and METEOR, at 0.5274 and 0.7310
(Tables~\ref{tab:finetune-baseline-epoch1-results} and
\ref{tab:finetune-low-lr-results}). The regularized-capacity configuration is
competitive for text realization, but is not the leader on either of these text
metrics or on triple F1.

The domain results are more informative than the averages alone. Fine-tuning
raises triple F1 markedly for GSMArena and Wikidata, while OpenWeather remains
close to zero for every configuration. The best reported OpenWeather F1 is only
0.0019 in the higher-capacity configuration. OWID also remains substantially
weaker than GSMArena and Wikidata. These differences are consistent with the
greater structural variation and compression required by weather forecasts and
statistical data, but the present experiments do not isolate which properties
of those domains cause the difficulty. They show that per-domain model
selection and additional domain-specific data work are needed rather than that
one hyperparameter setting transfers uniformly.

The fine-tuned model should also be understood as a distillation mechanism, not
as independent semantic supervision. Its targets are the reference texts and
normalized triples produced by the upstream pipeline. It can make the inferred
mapping from raw data to text and triples cheaper to execute in one decode pass,
but it is bounded by the semantic quality and coverage of those targets. Better
fine-tuning results therefore do not by themselves validate the underlying
triples against the original raw data. Improving target construction and adding
independently annotated examples remain necessary steps for establishing such
validity.

\section{Surface realization}

\subsection{Fitness function and domain effects}

The initial all-domain evolution run shows a clear domain difference. WrittenWork
has the highest BLEU score (0.5920), whereas Food has the lowest (0.0764), as
shown in Table~\ref{tab:surface-base}. WrittenWork may be easier because its
triples are often amenable to recurring bibliographic descriptions, whereas
food facts can require more varied lexical choices and combinations of
attributes. Nevertheless, the contrast justifies the subsequent
use of WrittenWork, Building, and Food as relatively easy, intermediate, and
difficult domains.

The three-domain variants demonstrate a trade-off between the selection signal
used during evolution and the metric used after evolution. The BLEU-fitness
variant without inspirations has the best BLEURT value among the reported
evolved programs, at 0.1049. Its paired variant with inspirations is slightly
higher on BLEU, 0.3526 versus 0.3521, but lower on BLEURT, 0.0952 versus 0.1049
(Tables~\ref{tab:surface-bleu-no-inspirations} and
\ref{tab:surface-bleu-with-inspirations}). Thus, removing inspirations cannot
be claimed to improve BLEU consistently; it improves BLEURT in this direct
comparison, while the BLEU difference is negligible and favours inspirations.

Variants optimized with LLM-as-a-judge fitness obtain stronger final judge
scores and lower detected additions and omissions than the BLEU-fitness
variants. For example, the no-inspiration LLM-judge variant has zero additions
and omissions and receives a Gemma judge score of 1.0000
(Table~\ref{tab:surface-llm-no-inspirations}). The mixed-model variants also
report zero additions and omissions. This pattern may be consistent with a partial 
alignment between the LLM-as-a-judge fitness criterion and the final semantic 
diagnostics: the fitness prompt rewards accurate reference to the input triples and 
penalizes extra information. However the zero scores 
cannot be attributed causally to LLM-judge fitness alone and do not 
constitute independent proof of semantic correctness. The
LLM-judge-with-inspirations variant still has nonzero additions, and the
reference-based BLEU, BLEURT, and METEOR values do not improve uniformly. The
selection objective therefore shapes the type of quality found by evolution;
it does not remove the need for external evaluation.

\subsection{Comparison with baselines}

The baseline comparison further emphasizes that there is no single dominant
surface-realization approach. The rule-based generator is clearly weaker than
the other systems across the reported measures. On the three selected domains,
the fine-tuned BART baseline has the best reference-based scores: BLEU 0.5628,
BLEURT 0.2398, and METEOR 0.7126. The zero-shot LLM prompt has lower reference-based 
scores, but it performs best in the final LLM-based assessments: the three 
overall quality judge scores from Gemma, Qwen, and Themis, and the separately prompted 
assessments of grammaticality, additions, and omissions. It receives the highest 
grammaticality score and has no detected additions or omissions.
(Table~\ref{tab:surface-final-comparison}). It is the strongest
reported system for judge-oriented faithfulness and fluency measures.

The reported evolved program does not yet exceed these stronger neural and
prompted baselines on the aggregate reference metrics. Its value lies elsewhere:
after evolution, the realization procedure is an inspectable Python artifact
that can be executed on unseen triple sets without another LLM call. This
property supports auditability, targeted modification, and potentially cheap
CPU execution. It does not, however, demonstrate a quality or cost advantage in
every deployment setting. A direct prompt is operationally simpler and is
currently more effective on several judged measures; BART remains stronger on
the principal reference metrics. The evolutionary approach is most compelling
where a reusable and inspectable realization program is a requirement, rather
than where the only objective is the highest benchmark score.

\section{Evaluation validity and limitations}

The judge--human comparison provides a necessary qualification to all
judge-based conclusions. For data-to-text, the LLM judge assigns a constant
score of 5 to every assessed instance, so Pearson, Spearman, and Kendall
correlations cannot be calculated. The corresponding quadratic weighted kappa
values are zero
(Table~\ref{tab:judge-human-text-correlation}). The automatic judge therefore
fails to distinguish the errors identified by the human evaluation in this
near-ceiling regime.

Text-to-triples agreement is more informative but remains mixed. For Evolving
Catalog omissions, positive associations are observed across the evaluated
domains; for example, the GSMArena omission correlation is 0.667 for Pearson,
Spearman, and Kendall. In contrast, Blueprint-Iterative Refinement has a
negative GSMArena additions correlation of $-0.509$, and several rows remain
undefined because one rater assigned a constant score
(Table~\ref{tab:judge-human-triples-correlation}).

Several aspects of the evaluation constrain generalization. The human study
contains ten instances per transition and domain and is scored by the thesis
author, so it is an exploratory expert evaluation rather than an
inter-annotator reliability study. The automatic pipeline evaluation covers
five domains, whereas the human evaluation and fine-tuning evaluation cover
four; ice hockey is not part of the human comparison. In addition, the use of
related LLM capabilities in both generation and judging can introduce shared
preferences. Reference metrics also assess similarity to a particular
lexicalization rather than all valid ways of expressing a triple set. Finally,
SENLEN is intentionally a simple period-delimited structural diagnostic and is
not a semantic evaluation measure. These limitations support combining metrics
and human inspection rather than treating any one score as definitive.

\section{Whole-system evaluation and future work}

The separate experiments establish the behavior of the two stages, but they do
not establish the quality of the complete system on raw data. The
whole-system evaluation, which was not done, addresses this gap by evolving a domain-specific
surface-realization program on the independent $S_3$ collection and evaluating
it on the separately collected $S_4$ collection. It would measure whether errors
or omissions in generated triples propagate into the final text, which cannot
be inferred reliably from isolated triple F1 or WebNLG scores.

The next evaluation cycle should therefore combine the completed end-to-end
test with a small independently annotated reference set, multiple human raters,
and agreement analysis between raters. Using judges from distinct model
families, calibrating them with contrastive error cases, and retaining
instance-level explanations would reduce the effect of score saturation.
Further surface-realization experiments should report all variants on the same
domains and include direct comparisons of inference latency, cost, and program
maintainability. Together, these steps would determine when the proposed
pipeline offers a practical advantage over direct prompting, rather than only
demonstrating that its components can be optimized separately.
